\documentclass[12pt]{article}

%%% Page setup %%%
\usepackage[top=1in, bottom=1in, left=1.2in, right=1in]{geometry}
\pagestyle{plain}
\linespread{1.5}

%%% Packages %%%
\usepackage[english]{babel}
\usepackage[utf8]{inputenc}
\usepackage{amsmath}
\usepackage[toc,page]{appendix}
\usepackage{fancyhdr}
\usepackage{gensymb}
\usepackage{graphicx}
\usepackage{listings}
\usepackage{parskip}
\usepackage{tabularx}
\usepackage{tikz}
\usepackage{wallpaper}
\usepackage{wrapfig}

\usepackage[hypertexnames=false]{hyperref}
\usepackage[numbered, depth=3]{bookmark}

\usepackage[backend=biber,style=apa]{biblatex}
\addbibresource{proposal.bib}

\usetikzlibrary{matrix}


%%% Variables %%%
\newcommand{\faculty}{Faculty of Technology}
\newcommand{\course}{Master's Thesis in Technology (DTEK1002)}
\newcommand{\documenttitle}{Research Proposal}
\newcommand{\authorname}{Florian Dejonckheere}
\newcommand{\authoremail}{fwdejo@utu.fi}
\newcommand{\documentdate}{\today}

%%% Document %%%
\begin{document}
	\begin{titlepage}
		\begin{center}
			\textsc{Turun Yliopisto \textemdash \hspace{1mm} University of Turku} \\
			\textsc{\faculty}
		\end{center}
		\begin{figure}[h]
			\vspace{10mm}

			\centering\includegraphics[width=0.8\textwidth]{logo-fi}

			\vspace{20mm}\label{fig:logo}
		\end{figure}
		\begin{center}
			\fontsize{10mm}{7mm}\selectfont
			\textup{\documenttitle} \\
			\textnormal{ \Large{\course}} \\

			\vspace{25mm}

			\large{\textbf{\authorname}} \\
			\large{\href{mailto:\authoremail}{\authoremail}} \\

			\vspace{35mm}

			\documentdate
		\end{center}
	\end{titlepage}

	\clearpage

	%%% Table of contents %%%

	\tableofcontents

	\cleardoublepage

	%%% Document body %%%

	\section{Abstract}\label{sec:abstract}

	Modern applications designed to run in the cloud are often programmed as a set of microservices: a decomposition of the application based on logical boundaries in the domain logic.
	This separation of concerns brings along a number of advantages such as scalability, fault isolation and independent deployment, over more traditional software development practices such as monolithic codebases.
	The adoption of a microservice-based architecture also introduces a number of challenges for the developers.
	It requires a more thorough understand of the domain logic, and increases the complexity of the codebase.

	In this thesis, we investigate and design a software architecture for deploying and managing monolithic codebases in dynamic languages as a set of microservices.
	The architecture transparently decomposes the application into separate microservices based on logical boundaries, and deploys them in a distributed fashion.
	Communication between the microservices is handled entirely by the architecture, so that the developers can concentrate on the domain logic of the application.
	The deployed application is scaled up and down based on the current load, and the microservices are migrated between nodes to optimize resource usage.
	This approach takes away the overhead of developing a microservice-based architecture, while retaining the advantages of such an architecture.

	The solution is designed to be language-agnostic, but we focus on the implementation for the Ruby programming language.
	We evaluate the architecture by implementing a proof-of-concept and comparing it to existing solutions.

	\clearpage

	\section{Aims and goals}\label{sec:aims-and-goals}

	\clearpage

	\section{Literature}\label{sec:literature}

	\clearpage

	\section{Research questions}\label{sec:research-questions}

	\clearpage

	\section{Research methods}\label{sec:research-methods}

	\clearpage

	\section{Contents}\label{sec:contents}

	\clearpage

	%%% Bibliography %%%

	\nocite{*}

	\printbibliography[heading=bibintoc]

\end{document}
